\QHeading{第14题　口算。}
\par\noindent


\begin{oralgrid}
\OralItem{1}{0.7+0.3}{1}
\OralItem{2}{1.2-0.5}{0.7}
\OralItem{3}{3.5+2.5}{6}
\OralItem{4}{1-0.4}{0.6}
\OralItem{5}{0.36+0.64}{1}
\OralItem{6}{4.6-0.6}{4}
\OralItem{7}{2.7+3.03}{5.73}
\OralItem{8}{5.8-1.8}{4}\end{oralgrid}
\par

\QHeading{口算思路补充（对应上表题号）}

\Answer{上面列出得数，下面按小题给出不同计算思路。计数单位法适用于小数，整数题可比较凑整、拆分或商不变等方法。}
\MethodTitle{第1小题：$0.7+0.3$}

\MethodBlock{口算思路一：按相同计数单位计算}{先看成若干个十分之一，按整数加减，最后还原成小数。}{\begin{aligned}[t]&\phantom{={}}0.7+0.3\\
&=(7+3)\div 10\\
&=10\div 10\\
&=1\end{aligned}}
\MethodTitle{第2小题：$1.2-0.5$}

\MethodBlock{口算思路一：按相同计数单位计算}{先看成若干个十分之一，按整数加减，最后还原成小数。}{\begin{aligned}[t]&\phantom{={}}1.2-0.5\\
&=(12-5)\div 10\\
&=7\div 10\\
&=0.7\end{aligned}}

\MethodBlock{口算思路二：补差法}{从减数先补到一个整数，再补到被减数，两段差相加。}{\begin{aligned}[t]&\phantom{={}}1.2-0.5\\
&=(1-0.5)+(1.2-1)\\
&=0.5+0.2\\
&=0.7\end{aligned}}
\MethodTitle{第3小题：$3.5+2.5$}

\MethodBlock{口算思路一：按相同计数单位计算}{先看成若干个十分之一，按整数加减，最后还原成小数。}{\begin{aligned}[t]&\phantom{={}}3.5+2.5\\
&=(35+25)\div 10\\
&=60\div 10\\
&=6\end{aligned}}
\MethodTitle{第4小题：$1-0.4$}

\MethodBlock{口算思路一：按相同计数单位计算}{先看成若干个十分之一，按整数加减，最后还原成小数。}{\begin{aligned}[t]&\phantom{={}}1-0.4\\
&=(10-4)\div 10\\
&=6\div 10\\
&=0.6\end{aligned}}
\MethodTitle{第5小题：$0.36+0.64$}

\MethodBlock{口算思路一：按相同计数单位计算}{先看成若干个百分之一，按整数加减，最后还原成小数。}{\begin{aligned}[t]&\phantom{={}}0.36+0.64\\
&=(36+64)\div 100\\
&=100\div 100\\
&=1\end{aligned}}
\MethodTitle{第6小题：$4.6-0.6$}

\MethodBlock{口算思路一：按相同计数单位计算}{先看成若干个十分之一，按整数加减，最后还原成小数。}{\begin{aligned}[t]&\phantom{={}}4.6-0.6\\
&=(46-6)\div 10\\
&=40\div 10\\
&=4\end{aligned}}

\MethodBlock{口算思路二：补差法}{从减数先补到一个整数，再补到被减数，两段差相加。}{\begin{aligned}[t]&\phantom{={}}4.6-0.6\\
&=(1-0.6)+(4.6-1)\\
&=0.4+3.6\\
&=4\end{aligned}}
\MethodTitle{第7小题：$2.7+3.03$}

\MethodBlock{口算思路一：按相同计数单位计算}{先看成若干个百分之一，按整数加减，最后还原成小数。}{\begin{aligned}[t]&\phantom{={}}2.7+3.03\\
&=(270+303)\div 100\\
&=573\div 100\\
&=5.73\end{aligned}}
\MethodTitle{第8小题：$5.8-1.8$}

\MethodBlock{口算思路一：按相同计数单位计算}{先看成若干个十分之一，按整数加减，最后还原成小数。}{\begin{aligned}[t]&\phantom{={}}5.8-1.8\\
&=(58-18)\div 10\\
&=40\div 10\\
&=4\end{aligned}}

\MethodBlock{口算思路二：补差法}{从减数先补到一个整数，再补到被减数，两段差相加。}{\begin{aligned}[t]&\phantom{={}}5.8-1.8\\
&=(2-1.8)+(5.8-2)\\
&=0.2+3.8\\
&=4\end{aligned}}


\par\noindent

\QHeading{第15题　用简便方法计算下面各题。}

\CalcBlock[方法一：简便计算]{1}{\begin{aligned}[t]&\phantom{={}}\color{black}15.9-2.48-7.52\\&\color{answerred}=15.9-(2.48+7.52)\\&\color{answerred}=15.9-10\\&\color{answerred}=5.9\end{aligned}}
\par
\MethodBlock{方法二：按运算顺序直接计算（对照）}{本法只作结果核对；本题要求简便计算，前面的简便解法仍保留。}{\begin{aligned}[t]&\phantom{={}}15.9-2.48-7.52\\
&=13.42-7.52\\
&=5.9\end{aligned}}


\CalcBlock[方法一：简便计算]{2}{\begin{aligned}[t]&\phantom{={}}\color{black}11.22-5.6+2.38\\&\color{answerred}=11.22+2.38-5.6\\&\color{answerred}=(11.22+2.38)-5.6\\&\color{answerred}=13.6-5.6\\&\color{answerred}=8\end{aligned}}
\MethodBlock{方法二：按运算顺序直接计算（对照）}{本法只作结果核对；本题要求简便计算，前面的简便解法仍保留。}{\begin{aligned}[t]&\phantom{={}}11.22-5.6+2.38\\
&=5.62+2.38\\
&=8\end{aligned}}

\CalcBlock[方法一：简便计算]{3}{\begin{aligned}[t]&\phantom{={}}\color{black}29.5-9.55+0.5\\&\color{answerred}=29.5+0.5-9.55\\&\color{answerred}=(29.5+0.5)-9.55\\&\color{answerred}=30-9.55\\&\color{answerred}=20.45\end{aligned}}
\MethodBlock{方法二：按运算顺序直接计算（对照）}{本法只作结果核对；本题要求简便计算，前面的简便解法仍保留。}{\begin{aligned}[t]&\phantom{={}}29.5-9.55+0.5\\
&=19.95+0.5\\
&=20.45\end{aligned}}

\CalcBlock[方法一：简便计算]{4}{\begin{aligned}[t]&\phantom{={}}\color{black}4.27+1.36+5.73\\&\color{answerred}=4.27+5.73+1.36\\&\color{answerred}=(4.27+5.73)+1.36\\&\color{answerred}=10+1.36\\&\color{answerred}=11.36\end{aligned}}
\MethodBlock{方法二：先处理括号或乘除，再凑整再减去多加的数}{把第二个数凑成整数，再减去多加的部分。}{\begin{aligned}[t]&\phantom{={}}4.27+1.36+5.73\\
&=5.63+5.73\\
&=5.63+6-0.27\\
&=11.63-0.27\\
&=11.36\end{aligned}}

\CalcBlock[方法一：简便计算]{5}{\begin{aligned}[t]&\phantom{={}}\color{black}13.45-2.8-7.2\\&\color{answerred}=13.45-(2.8+7.2)\\&\color{answerred}=13.45-10\\&\color{answerred}=3.45\end{aligned}}
\MethodBlock{方法二：按运算顺序直接计算（对照）}{本法只作结果核对；本题要求简便计算，前面的简便解法仍保留。}{\begin{aligned}[t]&\phantom{={}}13.45-2.8-7.2\\
&=10.65-7.2\\
&=3.45\end{aligned}}

\CalcBlock[方法一：简便计算]{6}{\begin{aligned}[t]&\phantom{={}}\color{black}6.18+2.57+3.43+3.82\\&\color{answerred}=6.18+3.82+2.57+3.43\\&\color{answerred}=(6.18+3.82)+(2.57+3.43)\\&\color{answerred}=10+6\\&\color{answerred}=16\end{aligned}}
\MethodBlock{方法二：交换、结合后先凑整}{保留每个数原来的加减号，先把能凑整的部分放在一起。}{\begin{aligned}[t]&\phantom{={}}6.18+2.57+3.43+3.82\\
&=(6.18+3.82)+2.57+3.43\\
&=10+2.57+3.43\\
&=12.57+3.43\\
&=16\end{aligned}}

\MethodBlock{方法三：按运算顺序直接计算（对照）}{本法只作结果核对；本题要求简便计算，前面的简便解法仍保留。}{\begin{aligned}[t]&\phantom{={}}6.18+2.57+3.43+3.82\\
&=8.75+3.43+3.82\\
&=12.18+3.82\\
&=16\end{aligned}}

\par\noindent

\QHeading{第16题　解决问题。}

\Problem{① 妈妈买了两只鸡，其中公鸡重2.5千克，公鸡比母鸡重600克。这两只鸡一共重多少千克？}
\par


\MethodTitle{方法一：简便计算}
\Answer{{\Formula{\begin{aligned}[t]&\phantom{={}}2.5+(2.5-0.6)\\&=2.5+2.5-0.6\\&=(2.5+2.5)-0.6\\&=5-0.6\\&=4.4\text{（千克）}\end{aligned}}}\qquad $600$克$=0.6$千克}
\Answer{答：两只鸡一共重4.4千克。}

\MethodTitle{方法二：先按两只都与公鸡同重计算}
\Answer{600克等于0.6千克。两只都按2.5千克算，多算了0.6千克。}
\Answer{\Formula{\begin{aligned}[t]&\phantom{={}}2.5\times 2-0.6\\
&=5-0.6\\
&=4.4\text{（千克）}\end{aligned}}}
\Answer{答：两只鸡一共重4.4千克。}
\Problem{② 一个笔袋2.5元，一本笔记本比一个笔袋贵1.6元。买一个笔袋和一本笔记本一共花多少元？}

\MethodTitle{方法一：简便计算}
\Answer{{\Formula{\begin{aligned}[t]&\phantom{={}}2.5+(2.5+1.6)\\&=(2.5+2.5)+1.6\\&=5+1.6\\&=6.6\text{（元）}\end{aligned}}}}
\Answer{答：一共花6.6元。}

\MethodTitle{方法二：先按两件同价计算，再补差价}
\Answer{两件都按2.5元计算，再加笔记本贵的1.6元。}
\Answer{\Formula{\begin{aligned}[t]&\phantom{={}}2.5\times 2+1.6\\
&=5+1.6\\
&=6.6\text{（元）}\end{aligned}}}
\Answer{答：一共花6.6元。}
\Problem{③ 小明每天骑行锻炼，第一天骑行25.45千米，第二天骑行24.55千米，第三天骑行30.8千米。小明三天一共骑行多少千米？}

\MethodTitle{方法一}
\Answer{{\Formula{\begin{aligned}[t]&\phantom{={}}25.45+24.55+30.8\\&=(25.45+24.55)+30.8\\&=50+30.8\\&=80.8\text{（千米）}\end{aligned}}}}
\Answer{答：三天一共骑行80.8千米。}

\MethodTitle{方法二：以每天25千米为基数进行补偿}
\Answer{第一天比25多0.45，第二天比25少0.45，正好抵消；第三天比25多5.8。}
\Answer{\Formula{\begin{aligned}[t]&\phantom{={}}25\times 3+5.8\\
&=75+5.8\\
&=80.8\text{（千米）}\end{aligned}}}
\Answer{答：三天一共骑行80.8千米。}
